% Part: sets-functions-relations
% Chapter: sets
% Section: russells-paradox
% Hindi translation (hi-Deva-IN) of Open Logic commit 9620cc73f9c8e0ad003c514a5d3748f29611c4c0.

\documentclass[../../../include/open-logic-section]{subfiles}


\begin{document}

\olfileid[hi]{sfr}{set}{rus}
\olsection{रसेल का विरोधाभास}

विस्तारिता $\Setabs{x}{\phi(x)}$ संकेतन की अनुमति देती है---ऐसे $x$ के
\emph{उस} समुच्चय के लिए जिनके लिए~$\phi(x)$ सत्य है। किंतु विस्तारिता
\emph{वास्तव में} केवल निम्न विचार की अनुमति देती है। \emph{यदि} कोई
ऐसा समुच्चय है जिसके अवयव ठीक वे ही वस्तुएँ हैं जो~$\phi$ को संतुष्ट
करती हैं, \emph{तो} ऐसा समुच्चय केवल एक है। दूसरे शब्दों में:
किसी~$\phi$ को नियत कर लेने पर समुच्चय
$\Setabs{x}{\phi(x)}$ अद्वितीय है, \emph{यदि उसका अस्तित्व हो}।

लेकिन यह सशर्त कथन महत्त्वपूर्ण है!{} निर्णायक बात यह है कि प्रत्येक
गुणधर्म \emph{समुच्चय-निर्माण} की अनुमति नहीं देता। अर्थात कुछ गुणधर्म
समुच्चय परिभाषित \emph{नहीं} करते। यदि सभी गुणधर्म ऐसा करते, तो सीधे
विरोधाभास उत्पन्न हो जाते। इसका सबसे प्रसिद्ध उदाहरण रसेल का
विरोधाभास है।

समुच्चय दूसरे समुच्चयों के !!{element} हो सकते हैं---उदाहरण के लिए,
समुच्चय~$A$ का घात समुच्चय स्वयं समुच्चयों से बना होता है। इसलिए यह
पूछना या जाँचना सार्थक है कि कोई समुच्चय दूसरे समुच्चय का
!!a{element} है या नहीं। क्या कोई समुच्चय स्वयं अपना अवयव हो सकता है?
समुच्चय की धारणा में ऐसा कुछ नहीं जो इसे असंभव ठहराए। उदाहरण के लिए,
यदि \emph{सभी} समुच्चय वस्तुओं का एक संग्रह बनाते हैं, तो यह मानना
स्वाभाविक लग सकता है कि उन्हें एक ही समुच्चय---सभी समुच्चयों के
समुच्चय---में एकत्र किया जा सकता है। और वह स्वयं समुच्चय होने के कारण
सभी समुच्चयों के समुच्चय का !!a{element} होगा।

रसेल का विरोधाभास तब उत्पन्न होता है जब हम स्वयं का !!a{element} न
होने के गुणधर्म, अर्थात \emph{स्वयं का अवयव न होने} पर विचार करते हैं।
यदि हम मान लें कि उन सभी समुच्चयों का एक समुच्चय है जो स्वयं अपने
!!a{element} नहीं हैं, तो क्या होगा? क्या
\[
R = \Setabs{x}{x \notin x}
\]
का अस्तित्व है? सिद्ध किया जा सकता है कि उसका अस्तित्व नहीं है।

\begin{thm}[रसेल का विरोधाभास]\ollabel{thm:russells-paradox}
	ऐसा कोई समुच्चय $R = \Setabs{x}{x \notin x}$ नहीं है।
\end{thm}

\begin{proof}
यदि $R = \Setabs{x}{x \notin x}$ का अस्तित्व हो, तो
$R \in R$ तभी और केवल तभी जब $R \notin R$; यह एक विरोधाभास है।
\end{proof}

\begin{tagblock}{novice}
\begin{explain}
आइए इस प्रमाण को कुछ धीरे-धीरे समझें। यदि $R$ का अस्तित्व है, तो यह
पूछना सार्थक है कि $R \in
R$ है या नहीं। मान लीजिए कि वास्तव में
$R \in R$। अब, $R$ को उन सभी समुच्चयों के समुच्चय के रूप में
परिभाषित किया गया था जो स्वयं अपने !!{element} नहीं हैं। इसलिए यदि
$R \in R$, तो $R$ में स्वयं $R$ का परिभाषक गुणधर्म नहीं है। किंतु
केवल वे समुच्चय ही~$R$ में हैं जिनमें यह गुणधर्म है; अतः $R$ स्वयं~$R$
का !!a{element} नहीं हो सकता, अर्थात $R \notin R$। लेकिन $R$ एक ही
समय में~$R$ का !!a{element} हो भी और न भी हो---यह संभव नहीं;
इसलिए हमें विरोधाभास मिलता है।

चूँकि $R \in R$ मानने से विरोधाभास मिलता है, अतः $R \notin R$। किंतु
इससे भी विरोधाभास मिलता है!{} यदि $R
\notin R$, तो $R$ में स्वयं
$R$ का परिभाषक गुणधर्म है; इसलिए अन्य सभी स्वयं का अवयव न होने वाले
समुच्चयों की तरह $R$ भी $R$ का !!a{element} होगा। फिर से, वह एक ही
समय में~$R$ का !!a{element} हो भी और न भी हो---यह संभव नहीं।
\end{explain}
\end{tagblock}

\begin{digress}
हम ऐसा समुच्चय सिद्धान्त कैसे बनाएँ जो रसेल के विरोधाभास में पड़ने से
बचे, अर्थात जो यह \emph{असंगत} दावा न करे कि
$R = \Setabs{x}{x \notin x}$ का अस्तित्व है? इसके लिए हमें ऐसे
अभिगृहीत निर्धारित करने होंगे जो अत्यंत सटीक शर्तें बताएँ कि समुच्चय कब
अस्तित्व में होते हैं (और कब नहीं होते)।
	
इस अध्याय में दी गई समुच्चय सिद्धान्त की रूपरेखा ऐसा नहीं करती। यह
\emph{सचमुच अनौपचारिक} है। यह केवल बताती है कि समुच्चय विस्तारिता का
पालन करते हैं और यदि कुछ समुच्चय दिए हों, तो उनका सम्मिलन, सर्वनिष्ठ
आदि बनाया जा सकता है। इससे अधिक कठोरता से समुच्चय सिद्धान्त विकसित करना
संभव है। \oliflabeldef{cumul:::part}{उस कठोर विवेचन को भाग
\olref[cumul][][]{part} के लिए सुरक्षित रखा जाएगा। अभी हम अनौपचारिक
रूप से आगे बढ़ेंगे और सावधानी से विरोधाभासों से बचने का प्रयास करेंगे।}{}
\end{digress}


\end{document}
